% Generated by scripts/assemble.py; corrections belong in source/corrections.json.
\examhead{159}{华南师范大学 2024 年数学分析试卷}{151}
% p152-b22; source PDF page 152
\par\Needspace{5\baselineskip}
\blockanchor{p152-b22}
\majorquestion{一}求极限. 每题 10 分, 共 50 分.\par
% p152-b23; source PDF page 152
\blockanchor{p152-b23}
\question{1}$\lim_{x\to 0}\frac{\sqrt{1 + \tan x} - \sqrt{1 - \tan x}}{\sin x}.$\par
% p152-b24; source PDF page 152
\blockanchor{p152-b24}
\question{2}$\lim_{x\to 1}\left(\frac{k}{1 - x^k} -\frac{l}{1 - x^l}\right)$ ，其中 $k,l$ 为正整数\par
% p153-b00; source PDF page 153
\blockanchor{p153-b00}
\question{3}$\lim_{n\to \infty}[\cos \pi \sqrt{n^2 + 1}]^2.$\par
% p153-b01; source PDF page 153
\blockanchor{p153-b01}
\question{4}$\lim_{n\to\infty}\frac{n}{\sqrt[n]{n!}}.$\par
% p153-b02; source PDF page 153
\blockanchor{p153-b02}
\question{5}$\lim_{\substack{x\to+\infty\\ y\to+\infty}}\left(\frac{xy}{x^{2}+y^{2}}\right)^{x^{2}}.$\par
% p153-b03; source PDF page 153
\par\Needspace{5\baselineskip}
\blockanchor{p153-b03}
\majorquestion{二}计算题. 每题 10 分, 共 40 分.\par
% p153-b04; source PDF page 153
\blockanchor{p153-b04}
\question{1}已知 $f(x)$ 在原点处存在二阶导数，且 $\lim_{x\to 0}\left(\frac{\sin x}{x^3} +\frac{f(x)}{x^2}\right) = 0,$ 求 $f'(0), f''(0).$\par
% p153-b05; source PDF page 153
\blockanchor{p153-b05}
\question{2}计算 $\int_{-2}^{-\sqrt{2}} \frac{1}{x \sqrt{x^2 - 1}} \, \mathrm{d}x$.\par
% p153-b06; source PDF page 153
\blockanchor{p153-b06}
\question{3}计算 $\iiint_{\Omega} x \, \mathrm{d}x \, \mathrm{d}y \, \mathrm{d}z, \Omega$ 是由平面 $x + 2y + z = 1, x = 0, y = 0, z = 0$ 围成的区域.\par
% p153-b07; source PDF page 153
\blockanchor{p153-b07}
\question{4}已知 $u$ 是由 $u = f(x, y, z, t), g(y, z, t) = 0, h(z, t) = 0$ 确定的关于 $x, y$ 的函数，求 $\frac{\partial u}{\partial x}, \frac{\partial u}{\partial y}$.\par
% p153-b08; source PDF page 153
\par\Needspace{5\baselineskip}
\par\Needspace{6\baselineskip}
\blockanchor{p153-b08}
\majorquestion{三}(10 分) 已知 $f(x)$ 是 $(0, +\infty)$ 上的单调递增函数, 且存在正数列 $\{a_n\}$ 满足\par
\nopagebreak[4]
% p153-b09; source PDF page 153
\blockanchor{p153-b09}
\displayformula{\lim _ {n \rightarrow \infty} a _ {n} = + \infty , \lim _ {n \rightarrow \infty} f (a _ {n}) = A.}
% p153-b10; source PDF page 153
\blockanchor{p153-b10}
证明： $\lim_{x\to +\infty}f(x) = A.$\par
% p153-b11; source PDF page 153
\par\Needspace{5\baselineskip}
\blockanchor{p153-b11}
\majorquestion{四}(10 分) 已知 $f(x)$ 在 $[a, b]$ 上连续且单调递增, 证明: $\int_{a}^{b} x f(x) \mathrm{d}x \geq \frac{a + b}{2} \int_{a}^{b} f(x) \mathrm{d}x$.\par
% p153-b12; source PDF page 153
\par\Needspace{5\baselineskip}
\blockanchor{p153-b12}
\majorquestion{五}（10分）证明函数项级数 $\sum_{n = 1}^{\infty}\frac{nx}{1 + n^5x^2}$ 在 $[0, + \infty)$ 上一致收敛.\par
% p153-b13; source PDF page 153
\par\Needspace{5\baselineskip}
\par\Needspace{6\baselineskip}
\blockanchor{p153-b13}
\majorquestion{六}(10 分) 已知 $f(x)$ 在 $[0, +\infty)$ 上连续, 无穷积分 $\int_0^{+\infty} f(x) \, \mathrm{d}x$ 收敛, 证明:\par
\nopagebreak[4]
% p153-b14; source PDF page 153
\blockanchor{p153-b14}
\displayformula{\lim _ {u \rightarrow 0 ^ {+}} \int_ {0} ^ {+ \infty} e ^ {- u x} f (x) \mathrm{d} x = \int_ {0} ^ {+ \infty} f (x) \mathrm{d} x.}
% p153-b15; source PDF page 153
\par\Needspace{5\baselineskip}
\par\Needspace{6\baselineskip}
\blockanchor{p153-b15}
\majorquestion{七}(20 分) 已知\par
\nopagebreak[4]
% p153-b16; source PDF page 153
\blockanchor{p153-b16}
\displayformula{f (x, y) = \left\{ \begin{array}{l l} (x ^ {2} + y ^ {2}) \sin \frac {1}{x ^ {2} + y ^ {2}}, & (x, y) \neq (0, 0); \\ 0, & (x, y) = (0, 0). \end{array} \right.}
% p153-b17; source PDF page 153
\blockanchor{p153-b17}
证明：\par
% p153-b18; source PDF page 153
\blockanchor{p153-b18}
(1) $f(x,y)$ 在 $(0,0)$ 处连续.\par
% p153-b19; source PDF page 153
\blockanchor{p153-b19}
(2) $f(x,y)$ 在 $(0,0)$ 处存在偏导数.\par
% p153-b20; source PDF page 153
\blockanchor{p153-b20}
(3) $f(x,y)$ 在 $(0,0)$ 处偏导数不连续.\par
% p153-b21; source PDF page 153
\blockanchor{p153-b21}
(4) $f(x,y)$ 在 $(0,0)$ 处可微.\par
